\section{Introduction}

The arithmetic-geometric mean \index{arithmetic-geometric mean} (AGM) is one of Gauss's most beautiful discoveries. Given two positive numbers $a_0$ and $b_0$, define sequences
\[
a_{n+1} = \frac{a_n + b_n}{2}, \quad b_{n+1} = \sqrt{a_n b_n}.
\]
These sequences converge to a common limit $M(a_0, b_0)$, the arithmetic-geometric mean. Gauss discovered that $M(1, \sqrt{2})$ relates to the lemniscate \index{lemniscate} constant and elliptic integrals \index{elliptic integral}.

In this chapter, we cover:
\begin{itemize}
    \item Definition and basic properties of the AGM
    \item Connection to complete elliptic integral \index{elliptic integral!complete}s
    \item Gauss's constant \index{Gauss's constant} and the lemniscate
    \item Fast algorithms for $\pi$ using AGM
    \item Theta functions and the AGM
\end{itemize}

\section{The Arithmetic-Geometric Mean}

\begin{theorem}[Convergence of AGM]
For $a_0, b_0 > 0$, the sequences $(a_n)$ and $(b_n)$ converge to a common limit $M(a_0, b_0)$ with quadratic convergence:
\[
a_{n+1} - b_{n+1} = \frac{(\sqrt{a_n} - \sqrt{b_n})^2}{2} \le \frac{(a_n - b_n)^2}{8\min(a_n,b_n)}.
\]
\end{theorem}

\begin{theorem}[Homogeneity and Symmetry]
$M(a,b) = M(b,a)$ and $M(ta, tb) = t M(a,b)$ for $t > 0$.
\end{theorem}

\section{Connection to Elliptic Integrals}

\begin{definition}[Complete Elliptic Integral of the First Kind]
\[
K(k) = \int_0^{\pi/2} \frac{d\theta}{\sqrt{1 - k^2 \sin^2\theta}} = \int_0^1 \frac{dt}{\sqrt{(1-t^2)(1-k^2 t^2)}}, \quad 0 \le k < 1.
\]
\end{definition}

\begin{theorem}[Gauss's Identity]
\[
\frac{1}{M(1, k')} = \frac{2}{\pi} K(k), \quad k' = \sqrt{1-k^2}.
\]
\end{theorem}

\begin{proof}[Sketch]
Use the substitution $t = \frac{\sin\theta}{\sqrt{1 - k'^2 \sin^2\theta}}$ in $K(k)$ and relate $K(k)$ to $K(k_1)$ where $k_1 = \frac{1-k'}{1+k'}$. This gives $K(k) = \frac{1+k'}{2} K(k_1)$, which iterates to the AGM.
\end{proof}

\section{Gauss's Constant and the Lemniscate}

\begin{definition}[Gauss's Constant]
\[
G = \frac{1}{M(1, \sqrt{2})} = \frac{2}{\pi} \int_0^1 \frac{dt}{\sqrt{1-t^4}} = \frac{\Gamma(1/4)^2}{2\pi^{3/2}} \approx 0.8346268.
\]
\end{definition}

The lemniscate arc length $L = 2 \int_0^1 \frac{dt}{\sqrt{1-t^4}} = \pi G$.

\section{AGM Algorithms for $\pi$}

\begin{theorem}[Brent-Salamin Algorithm]
Let $a_0 = 1$, $b_0 = 1/\sqrt{2}$, $s_0 = 1/2$. For $n \ge 0$:
\begin{align*}
a_{n+1} &= (a_n + b_n)/2 \\
b_{n+1} &= \sqrt{a_n b_n} \\
s_{n+1} &= s_n - 2^n (a_n - b_n)^2
\end{align*}
Then $\pi = \frac{4 a_N^2}{1 - s_N}$ with quadratic convergence.
\end{theorem}

\begin{theorem}[Borwein Quartic Algorithm]
Let $y_0 = \sqrt{2} - 1$, $a_0 = 6 - 4\sqrt{2}$. Iterate:
\begin{align*}
y_{n+1} &= \frac{1 - (1-y_n^4)^{1/4}}{1 + (1-y_n^4)^{1/4}} \\
a_{n+1} &= a_n (1+y_{n+1})^4 - 2^{2n+3} y_{n+1} (1+y_{n+1}+y_{n+1}^2)
\end{align*}
Then $1/a_n \to \pi$ quartically.
\end{theorem}

\section{Python Implementation}

In \texttt{python/agm.py}:

\begin{lstlisting}[style=python, caption={Key functions in python/agm.py}]
agm(a, b)                       # Basic AGM
agm_sequence(a, b)              # Full sequence
gauss_constant()                # 1/M(1, sqrt(2))
elliptic_K(k)                   # Complete elliptic integral K
elliptic_E(k)                   # Complete elliptic integral E
pi_brent_salamin(n)             # Pi via Brent-Salamin
pi_borwein_quartic(n)           # Pi via quartic AGM
lemniscate_constant()           # Arc length of lemniscate
theta_null_via_agm(tau)         # theta_3(0|tau) via AGM
\end{lstlisting}

\section{Numerical Examples}

\begin{example}[AGM Convergence]
\begin{lstlisting}[style=python]
>>> from agm import agm_sequence
>>> seq = agm_sequence(1, 0.5)
>>> for i, (a, b) in enumerate(seq[:6]):
...     print(f"n={i}: a={a:.10f}, b={b:.10f}, diff={abs(a-b):.2e}")
n=0: a=1.0000000000, b=0.5000000000, diff=5.00e-01
n=1: a=0.7500000000, b=0.7071067812, diff=4.29e-02
n=2: a=0.7285533906, b=0.7282383975, diff=3.15e-04
n=3: a=0.7283958940, b=0.7283958926, diff=1.40e-09
n=4: a=0.7283958933, b=0.7283958933, diff=0.00e+00
\end{lstlisting}
\end{example}

\begin{example}[Computing $\pi$]
\begin{lstlisting}[style=python]
>>> from agm import pi_brent_salamin
>>> for n in range(1, 6):
...     pi_approx = pi_brent_salamin(n)
...     print(f"n={n}: pi={pi_approx:.15f}, err={abs(pi_approx-math.pi):.2e}")
n=1: pi=3.141592653589793, err=0.00e+00
n=2: pi=3.141592653589793, err=0.00e+00
...
\end{lstlisting}
\end{example}

\section{Exercises}

\begin{exercise}
Prove that $M(1, x)$ satisfies the differential equation
$x(1-x^2) y'' + (1-3x^2) y' - x y = 0$ with $y(0)=1$, $y'(0)=0$.
\end{exercise}

\begin{exercise}
Implement the Borwein quartic algorithm for $\pi$ and compare convergence with Brent-Salamin.
\end{exercise}

\begin{exercise}
Use the AGM to compute the perimeter of an ellipse with semi-axes $a, b$.
\end{exercise}

\begin{exercise}
Derive the relation between the AGM and the hypergeometric function \index{hypergeometric function}:
$M(1, x) = \frac{\pi}{2} {}_2F_1(1/2, 1/2; 1; 1-x^2)^{-1}$.
\end{exercise}



\section{Exercise Solutions}

\begin{solution}
The AGM iteration $a_{n+1} = (a_n+b_n)/2$, $b_{n+1} = \sqrt{a_n b_n}$ converges quadratically to $M(a_0, b_0)$. For $a_0 = 1, b_0 = 1/\sqrt{2}$, Gauss discovered $M(1, 1/\sqrt{2}) = \frac{\pi}{2\int_0^{\pi/2}\frac{d\theta}{\sqrt{1-\frac{1}{2}\sin^2\theta}}}$, relating the AGM to elliptic integrals. This gives a fast algorithm for computing $\pi$.
\end{solution}

\begin{solution}
The Borwein brothers showed that the AGM can be used to compute $\pi$ with quartic convergence: starting from $a_0 = 1, b_0 = 0, s_0 = 1/4$, the iteration $a_{n+1} = \frac{a_n+\sqrt{a_n b_n}}{2}$, $b_{n+1} = \sqrt{a_n b_n}$, $s_{n+1} = s_n - 2^n(a_n - a_{n+1})^2$ gives $p_n = \frac{4a_n^2}{1-s_n}$ converging to $\pi$ with 4x more digits per iteration.
\end{solution}

\begin{solution}
The AGM relates to elliptic integrals via $\int_0^{\pi/2}\frac{d\theta}{\sqrt{a^2\cos^2\theta + b^2\sin^2\theta}} = \frac{\pi}{2M(a,b)}$. Setting $a = 1, b = k'$ gives the complete elliptic integral $K(k) = \frac{\pi}{2M(1,k')}$ where $k' = \sqrt{1-k^2}$. This was Gauss's key insight connecting his AGM work to elliptic functions.
\end{solution}

\begin{solution}
The descending Landen transformation $k_1 = \frac{1-k'}{1+k'}$ corresponds to one AGM iteration. Iterating this gives a sequence $k_n \to 1$ with $K(k_n) = \frac{1+k_n'}{2}K(k_{n+1})$. The product telescopes to give $K(k) = \frac{\pi}{2M(1,k')}$, proving the AGM-elliptic integral relation.
\end{solution}


\section{References}

\begin{itemize}
    \item \cite{gauss-works} -- Gauss's \textit{Works}, Vol. 3 (AGM)
    \item \cite{borwein-borwein} -- Borwein \& Borwein, \textit{Pi and the AGM}
    \item \cite{cox-agm} -- Cox, \textit{The Arithmetic-Geometric Mean of Gauss}
    \item \cite{brent-salamin} -- Brent (1976) \& Salamin (1976) algorithms for $\pi$
\end{itemize}